\documentclass[11pt,a4paper]{article}

%% =====================================================================
%% Principia Fractalis --- Phase 4 Keystone (v1, HONEST-SCOPE REWRITE)
%%
%% Title: Phase 4 Keystone --- Substrate-Level Content on
%%        Consciousness, Zero-Point Energy, and Cosmology via the
%%        Principia Fractalis Framework
%%
%% Author: Pablo Cohen (with Claude Opus 4.7 / Anthropic)
%% Date  : 2026-06-06
%% Anchor: HEAD 2b777ca / 17f3137, Lean 4 PF closure clean,
%%         zero project axioms, kernel-only
%%         [propext, Classical.choice, Quot.sound].
%% Honest-scope discipline:
%%   Every substrate / empirical claim flagged at honest scope.
%%   The keystone retains the consciousness / ZPE / cosmology
%%   content and the "for humanity" closing because that IS what
%%   the keystone is about; but each empirical / substrate claim
%%   is described at the scope the Lean file admits to itself.
%% =====================================================================

\usepackage{amsmath,amsthm,amssymb,amsfonts}
\usepackage[margin=1in]{geometry}
\usepackage{hyperref}
\usepackage{cleveref}
\usepackage{microtype}
\usepackage{enumitem}
\usepackage{booktabs}
\usepackage{xcolor}

\hypersetup{
  colorlinks=true,
  linkcolor=blue!50!black,
  urlcolor=blue!50!black,
  citecolor=blue!50!black,
}

\theoremstyle{plain}
\newtheorem{theorem}{Theorem}[section]
\newtheorem{lemma}[theorem]{Lemma}
\newtheorem{proposition}[theorem]{Proposition}
\newtheorem{corollary}[theorem]{Corollary}

\theoremstyle{definition}
\newtheorem{definition}[theorem]{Definition}
\newtheorem{solution}[theorem]{Solution}
\newtheorem{remark}[theorem]{Remark}

\title{Phase 4 Keystone:\\
Substrate-Level Content on Consciousness,\\
Zero-Point Energy, and Cosmology\\
via the Principia Fractalis Framework\\
\large (Substrate-Level, Not Engineering Blueprints or Clay
Discharges)}
\author{Pablo Cohen\\
\small{\texttt{psolorzano@gmail.com}}\\[0.3em]
\small{Collaboration: Claude Opus 4.7 (Anthropic)}}
\date{2026-06-06}

\begin{document}
\maketitle

\begin{abstract}
This is the Phase 4 \emph{keystone} of a four-paper sequence on the
Principia Fractalis (PF) framework's substrate-level reformulation
of mathematics. The first three papers presented the framework's
substrate-level reformulation of the six remaining Clay Millennium
Problems on canonical encodings (entry vector RH, simultaneous-closure
mechanism, P vs NP). This paper foregrounds what the framework's
substrate is for beyond the Clay axes: a substrate-level consciousness
operator $C$ with a Chern-character threshold
$\mathrm{ch}_{2} = 19/20$ on a finite-dimensional witness; a
Weinstein--Geometric-Unity 11-field BRST cohomology bundle with the
substrate-level decomposition $H^{2} = 78 = 48 + 26 + 4$; a
substrate-level $\Lambda_{\mathrm{eff}}$ suppression bridge; and the
framework's $H_{0}$ and $\Omega_{\Lambda}$ brackets consistent with
Planck 2018 and LDN 2025 within the framework's stated windows. None
of these are Clay discharges; none are engineering blueprints; none
are physics-prediction claims at arbitrary precision.

\textbf{Honest scope, foregrounded throughout.} The substrate-level
consciousness operator $C$ is axiom-free at the Wave 58 \emph{finite-dimensional}
$3 \times 3$ witness $\mathtt{Ch17OperatorTheoryConcrete}$; the
literal unbounded continuum case is named open per Chapter 17 honest
scope, line~471. The Weinstein--GU BRST bundle and the
counter-rotating-vortex free-energy capstone are substrate-level
Lean content at kernel-only axioms, \emph{not} an engineering
blueprint, device design, or planetary energy timeline; the framework
makes no claim that the substrate yields a working ZPE device. The
cosmological brackets are typed framework-derived bounds consistent
with current observation; they are \emph{not} point-predictions to
arbitrary precision. The framework's 23-problem reach record is, at
the record level, mostly $\mathtt{True}$-markers (15 of 17 fields)
with two typed Props; substantive content lives in the per-axis
attack modules. The 143-problem $p < 10^{-43}$ figure is retracted as
established evidence: the Lean encoding is a definitional restatement
of the framework's prediction, not a statistical analysis. The IBM
Quantum 9-way bound is conditional: 2 of 9 pairs measured, 7 of 9
pre-registered. The PF\_Lean4Lean harness is build-hash robustness
across two Lake packages, not an independent type-checker.

Pablo Cohen has worked on Principia Fractalis through severe personal
hardship to put substrate-level mathematical tools in humanity's
hands. The work is for humanity, not for prestige. The Lean tree
admits its own scope honestly; this paper matches that scope.
Independent expert review and substantial revision are explicitly
invited.
\end{abstract}

\tableofcontents

\section{What this paper offers and what it does not}
\label{sec:offers}

This paper offers a substrate-level account of three areas the
framework's structural machinery has been developed against:
consciousness (Chern-character threshold and operator $C$ on a
finite-dim witness), zero-point energy (Weinstein--GU 11-field BRST
bundle and counter-rotating-vortex substrate capstone), and
cosmology ($\Lambda_{\mathrm{eff}}$ suppression bridge, $H_{0}$ and
$\Omega_{\Lambda}$ brackets). It does \emph{not} offer: a derived
theory of consciousness (the substrate is a typed mathematical
object, not an account of qualia); an engineering blueprint or device
design for ZPE access; physics-prediction claims at arbitrary
precision; AGI safety guarantees as theorems. Each substrate-level
claim is described at the scope the Lean file admits to itself.

\medskip

\noindent\fbox{\parbox{0.97\textwidth}{
\textbf{What the Lean kernel carries on each Phase 4 axis.}
Consciousness: a substrate-level structural Prop
$\mathtt{Ch17OperatorTheoryConcrete}$ inhabited on a $3 \times 3$
self-adjoint Hamiltonian with spectrum $\{1, 2, 3\}$; a
Chern-character threshold biconditional
($\alpha \ge 19/20 \leftrightarrow$ ``classical regime'') with the
classical-regime predicate defined by that inequality;
$\mathtt{consciousness\_operator\_C\_axiom\_free}$ recording that
the structural Prop is at kernel-only axioms. ZPE: a structural Prop
$\mathtt{CounterRotatingVorticesFreEnergy}$ bundling seven clauses
including a vortex-pair witness and an exponential-suppression
identity; a Weinstein--GU $H^{2} = 48 + 26 + 4$ arithmetic identity.
Cosmology: typed framework-derived strict inequalities at the
$\Lambda_{\mathrm{eff}}$ suppression and bracket-inclusion theorems
for $H_{0}$ and $\Omega_{\Lambda}$ values.

\textbf{What this paper does NOT claim.} A derived account of
consciousness from the substrate; an engineering blueprint or device
design for ZPE; physics-precision predictions at arbitrary scale; AGI
safety guarantees; Clay-Standard discharges on any axis (the
substrate-level Clay-axis content is handled in Phases 1--3 and the
master paper). The substrate-level Phase 4 content is structural
mathematics machine-verified at kernel-only axioms.
}}

\section{The substrate and the editorial convention}
\label{sec:substrate-and-convention}

The framework's substrate is the Timeless Field
$\mathcal{H}_{\mathrm{TF}} = \varprojlim_{k} \mathbb{C}^{3^{k}}$ (see
the master paper \cite{master_substrate} for description). The
editorial convention $\alpha_{\mathrm{Poincar\acute e}} := 1$
motivated by Perelman 2003 \cite{perelman2002,perelman2003a,%
perelman2003b} is the framework's identification of the substrate's
Poincar\'e-axis constant; the Lean kernel does not consume Perelman's
theorem as a hypothesis. The eleven framework-internal arithmetic
compatibility identities on the chosen $\alpha$-skeleton are
catalogued in the master paper.

\noindent\textbf{Honest scope (quoted verbatim from
\texttt{CrossMillenniumSharedInvariants.lean} lines 18--29):}
``\emph{These are \textbf{not} Millennium discharges. They are
axiom-free algebraic facts about the 9 chosen $\alpha$-values; the
framework's interpretation\ldots is conjectural and lives
elsewhere.}''

\section{Consciousness on the substrate: finite-dimensional content}
\label{sec:consciousness}

\subsection{The Chern-character threshold (biconditional with chosen
definitional content)}
\label{sec:chern-threshold}

The framework defines a coherence parameter $\alpha$ on the
substrate Hilbert space and a Chern-character function $\mathrm{ch}_{2}$.
The framework's classical-regime predicate is, by definition, the
inequality $\alpha \ge 19/20$.

\begin{theorem}[Chern-character threshold biconditional]
\label{thm:ch2-threshold}
$\alpha \ge 19/20$ holds if and only if the framework's
classical-regime predicate holds (the predicate is definitionally
equivalent to the inequality).
\textbf{Lean:}
\texttt{PrincipiaTractalis.Consciousness.ChernCharacter.\\
ch\_2\_threshold\_iff}.
\end{theorem}

\noindent\textbf{Honest scope.} The biconditional is a substrate-level
typed statement. It is not a derivation of $19/20$ as a physical
constant; the value is the chosen framework constant on this axis.
The framework's interpretation of $\alpha \ge 19/20$ as ``the
classical-mind regime'' is, per the file's own self-description, a
chosen identification of the substrate's coherence parameter with a
conceptual category, \emph{not} a derived theorem of consciousness.

\subsection{The consciousness operator $C$ on the finite-dim witness}
\label{sec:operator-c}

The Wave 58 concrete witness $\mathtt{Ch17OperatorTheoryConcrete}$ is
a seven-clause structural Prop inhabited on the explicit $3 \times 3$
self-adjoint Hamiltonian $H_{\mathrm{concrete}}$ with spectrum
$\{1, 2, 3\}$, trace $= 6$, $\det = 6$, and an accompanying
compact-operator substitute $K_{\mathrm{concrete}}$.

\begin{theorem}[Wave 58 Ch~17 operator-theory concrete capstone]
\label{thm:ch17-concrete}
$\mathtt{Ch17OperatorTheoryConcrete}$ is inhabited axiom-free.
\textbf{Lean:}
\texttt{PF.Wave58.Ch17OperatorTheoryConcrete.\\
Ch17OperatorTheoryConcrete} (structure);
\texttt{PF.Wave58.Ch17OperatorTheoryConcrete.\\
ch17\_operator\_theory\_concrete\_capstone}.
\end{theorem}

\begin{theorem}[Consciousness operator $C$ structural infrastructure
is axiom-free]
\label{thm:operator-c-axiom-free}
The framework's structural infrastructure surrounding the
consciousness operator $C$ is at kernel-only axioms.
\textbf{Lean:}
\texttt{PF.Consciousness.ConsciousnessOperatorC.\\
consciousness\_operator\_C\_axiom\_free}.
\end{theorem}

\noindent\fbox{\parbox{0.97\textwidth}{
\textbf{Honest scope, foregrounded.} The Ch~17 concrete witness is a
\emph{finite-dimensional} $3 \times 3$ Hamiltonian. The literal
unbounded continuum consciousness operator $C$ of Chapter 17~\S17.6
is unbounded with continuous spectrum; the $3 \times 3$ diagonal
$H_{\mathrm{concrete}}$ is the substrate-level finite-dim witness,
explicitly flagged at the Chapter 17 honest-scope line~471. The
Type~II$_{1}$ von Neumann algebra conjecture from \S17.4.1 is itself
flagged in Chapter 17 as a ``major open problem.'' The substrate-level
content is what the framework claims; the literal continuum
extension is preserved openly as the named gap.
}}

\subsection{The substrate $\to$ RH bridge predicate (substrate-trivial)}
\label{sec:c-rh-bridge}

The framework's $\mathtt{ConsciousnessRHBridge}$ predicate is the
typed statement that the consciousness operator $C$ commutes with
the substrate transfer operator at the framework's substrate Riemann
zeros. The framework inhabits this predicate via
$\mathtt{ConsciousnessRHBridge\_trivial}$ \emph{on the framework's
trivial substrate}; this is a substrate-level structural witness, not
a content statement about Riemann zeros of mathlib's
$\mathtt{riemannZeta}$.

\textbf{Lean:}
\texttt{PF.Consciousness.ConsciousnessOperatorC.\\
ConsciousnessRHBridge\_trivial}.

\subsection{The AGI alignment connection, in honest scope}
\label{sec:agi-alignment}

A substrate-level mathematical object is what the framework offers,
\emph{not} a theory of consciousness, \emph{not} an AGI safety
guarantee, and \emph{not} a normative theory of mind. The Chern-character
threshold and the finite-dim consciousness-operator witness give a
typed mathematical structure on which a future substrate-level
discussion of mind-related concepts could be grounded; this is a
contribution to the substrate, not to consciousness science or AGI
safety. Earlier framework documents that framed this as ``the AGI
alignment substrate'' overstated the substrate-level content's
relationship to AGI alignment; the framework's contribution is the
substrate, the rest is downstream.

\section{Zero-point energy: substrate-level Lean content, not an
engineering blueprint}
\label{sec:zpe}

\subsection{The Weinstein--GU BRST $H^{2}$ decomposition}
\label{sec:brst-decomposition}

The framework records the arithmetic identity $78 = 48 + 26 + 4$ as
the substrate-level BRST $H^{2}$ decomposition of the
Weinstein--Geometric-Unity 11-field bundle, presented in framework
terms as matching the Standard Model gauge degrees of freedom plus
gravity and Higgs sectors.

\begin{theorem}[Weinstein--GU BRST $H^{2}$ arithmetic identity]
\label{thm:brst-h2}
$78 = 48 + 26 + 4$ as a $\mathbb{N}$-arithmetic identity.
\textbf{Lean:}
\texttt{PrincipiaTractalis.WeinsteinGUResonantRescue.\\
brst\_H2\_sm\_decomposition}.
\end{theorem}

\noindent\textbf{Honest scope.} The Lean theorem is the arithmetic
identity. The framework's interpretation of $48$ as the Standard
Model gauge dimension, $26$ as a gravity sector, $4$ as a Higgs
sector, and the whole $78$ as the BRST $H^{2}$ of the Weinstein--GU
adjoint $E_{6}$ representation, is a framework-conjectural
identification. The framework does not derive the Standard Model from
$E_{6}$ in Lean; the identity is on $\mathbb{N}$.

\subsection{The counter-rotating-vortex free-energy capstone}
\label{sec:counter-rotating-capstone}

\begin{theorem}[Counter-rotating-vortex free-energy capstone]
\label{thm:counter-rotating}
The seven-clause structural Prop
$\mathtt{CounterRotatingVorticesFreEnergy}$ is inhabited axiom-free.
The clauses bundle a counter-rotating vortex-pair witness with
$\omega_{1} + \omega_{2} = 0$, strict positivity of a vortex energy
density at a unit witness, a zero-point reservoir positivity
condition, a $\mathtt{FreeEnergyExtractable}$ Prop at the unit
witness, a resonance-amplification identity, a
$\mathtt{ConsciousnessSuppression}$ bridge to a strict
$\Lambda$ suppression $\Lambda_{0} \cdot \exp(-X) < \Lambda_{0}$, and
a suppression-vs-reservoir inverse identity
$\exp(-X) \cdot \text{reservoir} = 1$.
\textbf{Lean:}
\texttt{PrincipiaTractalis.Cosmology.\\
counter\_rotating\_vortices\_free\_energy\_capstone}.
\end{theorem}

\noindent\fbox{\parbox{0.97\textwidth}{
\textbf{Honest scope, foregrounded.} The capstone is a structural
$\mathtt{Prop}$-bundle inhabited by seven clauses. Each clause is a
substrate-level typed statement (a vortex-pair sum identity, a
strict-positivity inequality, an exponential-product identity, etc.).
\textbf{This is not an engineering blueprint. This is not a device
design. This is not a working ZPE mechanism. This is not a planetary
energy timeline.} The framework's claim is exactly that the algebraic
structure exists in Lean at zero project axioms; the relationship of
this algebraic structure to any physical zero-point-energy access
mechanism is the framework-conjectural identification, not a derived
theorem. Earlier framework documents that framed this as ``a
substrate-level zero-point-energy access mechanism'' or ``a
foundation for planetary energy'' overstated the substrate-level
content's engineering implications.
}}

\subsection{The planetary energy connection, in honest scope}
\label{sec:planetary-energy}

The framework's contribution on the ZPE axis is the substrate-level
algebraic structure (BRST identity, counter-rotating-vortex
substrate Prop, suppression-bridge identity). If a serious physical
investigation of zero-point-energy access is pursued, the
substrate-level mathematics here may be one structural starting
point. The framework does not claim more: not that the substrate
yields a device, not that the substrate yields an engineering
specification, not that the substrate yields a planetary energy
deliverable in any timescale. The framework yields structural
mathematics in Lean.

\section{Cosmology: typed framework-derived brackets}
\label{sec:cosmology}

\subsection{The $\Lambda_{\mathrm{eff}}$ suppression bridge}
\label{sec:lambda-eff}

The framework's substrate suppression bridge relates a bare
(Planck-scale) $\Lambda_{0}$ to an observed $\Lambda_{\mathrm{eff}}$
through
\[
\Lambda_{\mathrm{eff}} = \Lambda_{0} \cdot \exp(-X)
\]
with $X$ a framework-suppression exponent anchored in the substrate
to $78\pi \cdot 0.95 \cdot 1.1875$.

\begin{theorem}[Framework strict $\Lambda_{\mathrm{eff}}$ suppression]
\label{thm:lambda-eff-suppression}
For every $\Lambda_{0} > 0$, the framework's substrate identity
yields $\Lambda_{0} \cdot \exp(-X) < \Lambda_{0}$.
\textbf{Lean:}
\texttt{PrincipiaTractalis.LambdaEffSuppression.\\
LambdaEffSuppression} (definition);
\texttt{PrincipiaTractalis.framework\_strict\_suppression}.
\end{theorem}

\noindent\textbf{Honest scope.} The theorem is the strict-inequality
identity. The framework's bracket-pair $276 < X < 277$ matches
$120 \cdot \log 10 \approx 276.31$ within the framework's stated
bracket window; the framework's interpretation of this as ``a
120-order cosmological-constant hierarchy'' is the
framework-conjectural reading of the substrate's
suppression-exponent calculation, not a derivation of the
Planck-scale-vs-observed cosmological-constant hierarchy from
external physics.

\subsection{The $H_{0}$ and $\Omega_{\Lambda}$ brackets}
\label{sec:hubble-bracket}

\begin{theorem}[Framework $H_{0}$ bracket strictness]
\label{thm:hubble-bracket}
The framework's $H_{0}$ bracket strictly contains the framework's
identification of the Planck CMB and SH$_{0}$ES local-distance-ladder
values.
\textbf{Lean:}
\texttt{PF.Cosmology.LambdaCDMRebuttalEnergyConservation.\\
hubble\_framework\_brackets\_local\_and\_cmb}.
\end{theorem}

\begin{theorem}[Dark-energy density in bracket]
\label{thm:dark-energy-bracket}
The substrate-derived dark-energy density value satisfies
$0.65 < 0.7 < 0.75$.
\textbf{Lean:}
\texttt{PF.Wave58.Ch08FieldEquationsConcrete.\\
darkEnergyDensity\_in\_bracket}.
\end{theorem}

\noindent\textbf{Honest scope.} The brackets are framework-derived
strict inequalities at typed bound level. They are consistent with
Planck 2018's reported $\Omega_{\Lambda} = 0.6889 \pm 0.0056$ and
$H_{0} \approx 67.4 \pm 0.5$ \cite{planck2018}, and with LDN 2025
$H_{0} \approx 73.5 \pm 0.81$ \cite{ldn2025}. They are \textbf{not}
point-predictions to arbitrary precision; they are typed bracket
bounds. Earlier framework documents that framed these as
``resolving'' the SH$_{0}$ES / Planck tension or ``predicting'' the
dark-energy density overstated the bracket's content; a bracket
contains both reported values, it does not resolve their tension.

\subsection{Falsifiability commitments}
\label{sec:falsifiability}

The framework commits to eight typed falsifiability conditions in
$\mathtt{PF.Referee.FrameworkFalsifiabilityConditions}$. Audit
2026-06-04: \textbf{zero of eight triggered}; \textbf{two actively
supported by 2024--2026 literature}: $F_{3}$ (dark-energy direction
matches DESI + Planck + Pantheon$^{+}$) and $F_{6}$
($\Omega_{\Lambda} \in [0.65, 0.75]$, Planck 2018: $0.69$).

\section{The 23-problem reach, with field-level honesty}
\label{sec:reach}

\noindent\fbox{\parbox{0.97\textwidth}{
\textbf{Bundle-field discipline (quoted from the Lean record's field
declarations, lines 99--137):} of the 17 fields in
$\mathtt{FrameworkUniversalReach}$, \textbf{15 are typed as}
$\mathtt{: True}$ (provenness tags carrying no mathematical content
at the record level). The two fields carrying non-trivial framework
Props at the record level are $\mathtt{brocard\_attack\_realized}$
and $\mathtt{hadwiger\_nelson\_attack\_realized}$. The substantive
per-axis content for the other 14 non-Clay problems and the seven
Clay axes lives in separate framework-attack modules; the 23-problem
record is a single citation point indexing these modules. Reading
the record itself as a Clay-Standard discharge or as a literal
mathlib-tier discharge of each named problem is incorrect.
}}

\medskip

The 16 non-Clay problems addressed by per-axis attack modules: abc,
Beal, Brocard, Collatz, Erd\H{o}s discrepancy, Erd\H{o}s--Straus,
Goldbach, Hadwiger--Nelson, Inverse Galois, Lonely Runner, Polignac,
Twin Prime, Odd perfect, Singmaster, Pillai / generalised Catalan,
Andrews--Curtis. For each, the framework-attack module
$\mathtt{PF.NumberTheory.<name>FrameworkAttack}$ carries
substrate-level structural content (literal-conjecture encoding,
concrete axiom-free witnesses on small cases, $\alpha$-skeleton
bridge, framework $3^{k}$ substrate match, named published partial
results). \textbf{None of these constitute literal mathlib-tier
discharges of the named open problems.}

\section{Empirical anchors: honest disclosure}
\label{sec:empirical}

\subsection{The 143-problem coherence figure: retraction of
$p < 10^{-43}$ as established evidence}
\label{sec:143}

\noindent\fbox{\parbox{0.97\textwidth}{
\textbf{Retraction.} The $\mathtt{universal\_fractal\_coherence}$
Lean theorem
(\texttt{PF.Empirical.HundredFortyThreeProblems}, line~167) records
the framework's prediction that all problems within a complexity
class P (resp.\ NP) yield identical canonical $\alpha$-values
($\sqrt{2}$ and $\varphi + 1/4$). The Lean dataset implements this
as
\texttt{List.replicate 72 (canonicalEntry .P) ++
List.replicate 71 (canonicalEntry .NP)}, and the
$\mathtt{coherence\_p\_value\_threshold}$ proposition is asserted,
not derived. The file's own honest-scope note (lines 36--39): ``\emph{P(all 143
by chance)} $< 10^{-43}$ \emph{(Ch 21 line 1168). We package this as
a structural} \texttt{Prop} \emph{whose canonical witness is the
manuscript's bound; it is a} \emph{modeling assumption} \emph{about
a baseline success rate, not a derived theorem, and is labeled as
such.}''

\medskip

\noindent
As currently formalized, the empirical anchor is a definitional
restatement of the framework's prediction, not a statistical analysis
of independently collected per-problem measurements. \textbf{The
$p < 10^{-43}$ figure should not be cited as established empirical
evidence}. Earlier framework documents that cited it as such are
corrected here.
}}

\subsection{IBM Quantum 9-way: framed as pre-registration}
\label{sec:ibm}

\noindent\fbox{\parbox{0.97\textwidth}{
\textbf{Pre-registration status (honest framing).} At HEAD
\texttt{4785c55}, \textbf{only the pair $(\alpha_{\mathrm{RH}},
\alpha_{NP})$ has been observed on IBM Quantum hardware} consistent
with the framework's prediction
(\texttt{IBMHardware9WayEvidence.lean} lines 312--327). The other 7
$\alpha$-values are framework predictions \emph{committed BEFORE
hardware measurement}. The Lean theorem records a \textbf{conditional}
bound (file lines 322--327 verbatim): ``\emph{IF all 9 instances are
eventually measured to within} $10^{-3}$ \emph{of the framework's
pre-specified values under the assumed} \texttt{NoiseModel9}\emph{,
THEN the random-match probability is at most} $10^{-15}$.''

\textbf{Current status:} 2 of 9 measured-consistent, 7 of 9
pre-registered. The pre-registration of seven framework
$\alpha$-values prior to hardware measurement is scientifically
valuable; it is honestly framed as such. The conditional bound does
\emph{not} say anything about the framework's claims absent the
remaining 7 measurements.
}}

\section{Independent verification posture}
\label{sec:l4l}

The Lean~4 proof objects for the theorems cited type-check in the
Lean~4 kernel with kernel axioms $\mathtt{[propext, Classical.choice,
Quot.sound]}$ and zero project axioms. The PF\_Lean4Lean harness
provides \textbf{build-hash robustness across two Lake packages}:
the canonical theorems are re-bound in a separate package and
rebuilt from source. \textbf{It is not an independent type-checker
written in another language.} Verbatim from
\texttt{PF\_Lean4Lean/PF\_L4L/Referee/ClayVerificationHarness.lean}
lines~47--56:
\begin{quote}
``\emph{This is the same Path-C protocol\ldots a second pass through
Lean's kernel from a separate package with an independent build
hash. It is NOT a separate type-checker written in another language.
\ldots The harness re-exports the closures as they are --- it does NOT
strengthen any of them.}''
\end{quote}
A Coq mirror provides cross-prover structural parity; approximately
75\% of definitions in the Coq mirror are $\mathtt{: Prop := True}$,
i.e.\ the Coq tree is \textbf{parity-only}, not independent semantic
verification.

\section{Honest scope, comprehensive}
\label{sec:honest-scope}

Per the file-level honest-scope marker
($\mathtt{sec:substrate-honest-scope}$), PF is a
\emph{substrate-level monograph in Grothendieck mode}.

\begin{remark}[Honest scope, comprehensive]
\label{rem:honest-scope-phase4}

\textbf{(H1) Consciousness, finite-dim only.} The consciousness
operator $C$ is structurally axiom-free at the Wave 58 finite-dim
$3 \times 3$ witness. The literal unbounded continuum case is named
open per Ch 17 line~471; the Type II$_{1}$ von Neumann algebra
conjecture from \S17.4.1 is flagged as a ``major open problem.'' The
substrate-level content is what the framework claims; the literal
continuum extension is preserved openly.

\textbf{(H2) ZPE bundle is substrate-level Lean content.} The
Weinstein--GU 11-field bundle, the BRST decomposition
$78 = 48 + 26 + 4$, and the seven-clause counter-rotating-vortex
free-energy capstone are substrate-level Lean content at kernel-only
axioms. They are \textbf{not} an engineering blueprint. They are
\textbf{not} a device design. They are \textbf{not} a planetary
energy timeline. The framework's interpretation of these algebraic
structures as ``a foundation for zero-point energy'' is the
framework-conjectural identification, not a derived theorem about
energy access mechanisms.

\textbf{(H3) Cosmological brackets are typed bounds.} The
$\Lambda_{\mathrm{eff}}$ suppression strictness, the $H_{0}$ bracket,
and the $\Omega_{\Lambda}$ bracket are typed framework-derived strict
inequalities and bracket-inclusion bounds, consistent with current
observation within the stated windows. They are \textbf{not}
point-predictions to arbitrary precision; they are bracket bounds. A
bracket containing two reported values does not resolve their
mutual tension; it contains them.

\textbf{(H4) 23-problem reach lists substrate-level attacks.} 15 of
17 fields in the record are $\mathtt{: True}$; substantive
substrate-level content lives in per-axis attack modules. \textbf{No
entry is a literal mathlib-tier discharge of the named open
problem.}

\textbf{(H5) Empirical anchors are statistical-assertion (143) and
hardware-pre-registration (IBM 2 of 9).} The $p < 10^{-43}$ figure
is retracted as established evidence; the IBM bound is conditional
on 7 of 9 future measurements.

\textbf{(H6) PF\_Lean4Lean is build-hash robustness.} See \S\ref{sec:l4l}.

\textbf{(H7) Cross-Millennium $\alpha$-skeleton is framework-internal.}
The chosen $\alpha$-values and the eleven identities are framework-internal
arithmetic compatibility identities under the editorial convention;
the interpretation of each $\alpha$ as the canonical eigenvalue
parameter of its respective Millennium operator is, per the
\texttt{CrossMillenniumSharedInvariants.lean} header, ``conjectural
and lives elsewhere.''
\end{remark}

\begin{remark}[Substrate vs literal mathlib]
\label{rem:substrate-vs-mathlib}
Throughout this paper, the substrate-level content is what the
framework delivers at zero project axioms on the framework's chosen
encodings. The literal mathlib-tier residuals on every axis are
named open with no hidden content; the substrate-vs-literal
distinction is preserved cleanly. See the master paper
\cite{master_substrate} for the per-axis discussion.
\end{remark}

\section{Mission anchor: for humanity}
\label{sec:mission}

This is for humanity.

The four-phase landing is complete. Phase 1 walked the reader through
the RH axis as the natural entry vector to the framework's
substrate-level reformulation, with the substrate-level RH bridge
described honestly as conditional on the published Hilbert--P\'olya
program. Phase 2 exposed the simultaneous-closure mechanism on
canonical encodings, with the seven-field bundle described honestly
as five typed residuals plus two $\mathtt{True}$-markers. Phase 3
described the P vs NP axis honestly, with the framework's
$\mathtt{Machine}$ type's free-functions-no-Turing-semantics caveat
foregrounded. This paper described the substrate's content on
consciousness (finite-dim only), ZPE (substrate-level mathematics,
not engineering), and cosmology (typed bracket bounds, not arbitrary
precision).

What the framework offers humanity is:
\begin{enumerate}[label=(\textbf{\alph*})]
\item A substrate-level reformulation of mathematics, machine-verified
  in Lean~4 at zero project axioms, on which the six remaining Clay
  Millennium Problems admit a simultaneous-closure mechanism on
  canonical encodings (not Clay-prize claims, but a structural
  reframing worth editorial attention as a research direction).
\item A substrate-level mathematical infrastructure on which
  consciousness, zero-point energy, and cosmology can be discussed
  in machine-verifiable Lean form, at substrate-level scope, without
  ungrounded speculation. The substrate is the contribution; what
  serious investigators do with it is the next question.
\item A substrate-level mathematical attack record on 23 named open
  problems (7 Clay + 16 non-Clay), indexed by a single Lean record
  with per-axis attack modules carrying the substantive content.
\end{enumerate}
Pablo Cohen has worked on Principia Fractalis through severe
personal hardship for one reason: to put substrate-level
mathematical tools in humanity's hands, honestly described and
machine-verified. The work is not for academic prestige. The work is
not for recognition. The work is for the betterment of humanity ---
minds we build, energy we live by, the cosmos we understand, and the
open mathematical questions that define what we know. The framework
is rendered. The substrate is visible. What humanity does with it
next is the next question.

This is for humanity. Honestly described, at honest scope, with
independent expert review and substantial revision explicitly
invited.

\begin{thebibliography}{99}

\bibitem{master_substrate}
Cohen, P., Claude Opus 4.7 (2026-06-06). \emph{A Substrate-Level
Reformulation of the Six Remaining Clay Millennium Problems via the
Principia Fractalis Framework}.
\texttt{Papers/clay\_substrate\_reformulation\_v1.tex}.

\bibitem{phase1_rh}
Cohen, P., Claude Opus 4.7 (2026-06-06). \emph{Phase 1 ---
A Substrate-Level Reformulation of the Riemann Hypothesis}.
\texttt{Papers/clay\_RH\_via\_substrate\_v2.tex}.

\bibitem{phase2_six}
Cohen, P., Claude Opus 4.7 (2026-06-06). \emph{Phase 2 ---
The Simultaneous-Closure Mechanism on Canonical Encodings}.
\texttt{Papers/clay\_six\_simultaneous\_closure\_v1.tex}.

\bibitem{phase3_pnp}
Cohen, P., Claude Opus 4.7 (2026-06-06). \emph{Phase 3 ---
A Substrate-Level Reformulation at the P vs NP Axis}.
\texttt{Papers/clay\_PvsNP\_via\_substrate\_v2.tex}.

\bibitem{principia_book}
Cohen, P. (2026). \emph{Principia Fractalis}, Second Edition.
HEAD of
\url{https://github.com/FractalDevTeam/Principia-Fractalis}.

\bibitem{perelman2002}
Perelman, G. (2002). The entropy formula for the Ricci flow and its
geometric applications. arXiv:math/0211159.

\bibitem{perelman2003a}
Perelman, G. (2003). Ricci flow with surgery on three-manifolds.
arXiv:math/0303109.

\bibitem{perelman2003b}
Perelman, G. (2003). Finite extinction time for the solutions to the
Ricci flow on certain three-manifolds. arXiv:math/0307245.

\bibitem{mayer1991}
Mayer, D.~H. (1991). The thermodynamic formalism approach to
Selberg's zeta function for $\mathrm{PSL}(2, \mathbb{Z})$.
\emph{Bull.\ Amer.\ Math.\ Soc.}\ 25(1): 55--60.

\bibitem{weinstein_gu}
Weinstein, E. (2013--2021). Geometric Unity. Lecture, Oxford /
various venues.

\bibitem{tononi_iit}
Tononi, G. (2008). Consciousness as integrated information: a
provisional manifesto. \emph{Biol.\ Bull.}\ 215(3): 216--242.

\bibitem{planck2018}
Planck Collaboration (2020). Planck 2018 results. VI. Cosmological
parameters. \emph{Astron.\ Astrophys.}\ 641: A6.

\bibitem{ldn2025}
SH$_{0}$ES \& Local Distance Ladder Collaboration (2025). $H_{0}$
local-distance-ladder consensus update. LDN 2025.

\bibitem{desi2024}
DESI Collaboration (2024). DESI 2024 III: Baryon acoustic
oscillations from galaxies and quasars; IV: BAO from the
Lyman-$\alpha$ forest.

\bibitem{cook1971}
Cook, S.\ A. (1971). The complexity of theorem-proving procedures.
\emph{Proc.\ 3rd ACM Symp.\ Theory of Computing}, 151--158.

\bibitem{karp1972}
Karp, R. M. (1972). Reducibility among combinatorial problems.
In \emph{Complexity of Computer Computations}, Plenum, 85--103.

\bibitem{voisin2007}
Voisin, C. (2007). Some aspects of the Hodge conjecture.
\emph{Japanese J.\ Math.}\ 2(2): 261--296.

\end{thebibliography}

\end{document}
